% problem-000383 1 分析化学定量方法
\section*{1. 分析化学定量方法}
\textit{下列为分问。}

\subsection*{1-1}
同族元素D， $\mathrm { { \bf E } } , \mathrm { { \bf ~ F } }$ （均为⾮放射性副族元素）基态原⼦的价层电⼦组态符合同⼀个通式，在元素周期表中，位置在A、B、C所在族之前，请给出D、E、F的元素符号及价层电⼦组态的通式。

\subsection*{1-2}
元素A存在于动物的⾎蓝蛋⽩中，⼈对A元素代谢紊乱可导致Wilson病。

$\mathbf { 1 - 2 - 1 \ A ^ { 2 + } }$ 硫酸盐在碱性溶液中加⼊缩⼆脲 $\mathrm { { H N } ( C O N H _ { 2 } ) _ { 2 } }$ 会得到特征的紫⾊物质，该物质为−2价的配位阴离⼦，具有对称中⼼和不通过A的⼆重旋转轴，⽆⾦属-⾦属键，请画出该阴离⼦的结构（A必须写元素符号，下同）。

\subsection*{1-2}
-2 通过 $\mathbf { A } ^ { 2 + }$ 与过量的丁⼆酮肟形成⼆聚配合物，实现了 $\mathbf { A } ^ { 2 + }$ 的平⾯正⽅形配位向 $\mathbf { A } ^ { 2 + }$ 的四⽅锥形配位转化，请画出该配合物的结构（丁⼆酮肟⽤ $\widehat { \mathrm { ~ N ~ } } \widehat { \mathrm { ~ N ~ } }$ 表⽰）。

\subsection*{1-3}
⽤ $\mathbf { B } ^ { + }$ 的标准溶液滴定KCl和KSCN的中性溶液，得到电位滴定曲线，其拐点依次位于M、N、P处。1-3-1 请分别写出在M、N、P处达到滴定终点的离⼦反应⽅程式。

\subsection*{1-3}
-2 在N处的物质是⽆⽀链的聚合物，请⾄少画出三个单元表⽰其结构。

\subsection*{1-4}
元素C的单质与单⼀的⽆机酸不起作⽤，但可溶于王⽔。C与 $\mathrm { O _ { 2 } }$ 和 $\mathrm { F } _ { 2 }$ 作⽤制得化合物X，在X中C的质量分数为57.43%，X的结构与1962年N. Bartlett开创性⼯作的产物极为相似，通过推演给出X的化学式。
